\documentclass[aspectratio=169,10pt,english]{beamer}

\input{../../../_preambulo_beamer.tex}
\input{../../../_comandos_beamer.tex}

\selectlanguage{english}

\title{MA1001B - Statistical Modeling for Decision Making}
\subtitle{Lesson 07: Counting Techniques II}
\author{}
\institute{Tecnol\'ogico de Monterrey}
\date{}

\begin{document}

\begin{frame}[plain]
  \vspace{1.2cm}
  {\Large\bfseries MA1001B - Statistical Modeling for Decision Making\par}
  \vspace{0.7cm}
  {\large Lesson 07: Repeated Categories and Allocations\par}
  \vspace{0.7cm}
  {\color{TecRojo}\rule{\textwidth}{1.2pt}}
  \vspace{0.8cm}

  MA1001B Faculty\\[0.3cm]
  Term\\[0.3cm]
  Tecnol\'ogico de Monterrey
\end{frame}

\begin{frame}{The Allocation Question}
  \begin{alertblock}{Guiding question}
    How do repeated categories change a count, and when does a correct count
    still fail to determine a probability?
  \end{alertblock}

  \vspace{0.4em}
  By the end of this lesson, you can:
  \begin{itemize}
    \item count sequences with repeated category labels;
    \item interpret multinomial coefficients as fixed-size allocations;
    \item count nonnegative allocations with stars and bars;
    \item explain why possible profiles need not be equally likely.
  \end{itemize}

  \vfill
  \scriptsize All requests, queues, centers, and routing mechanisms are synthetic.
\end{frame}

\begin{frame}{Repeated Permutations Remove Duplicate Rearrangements}
  \small
  Eight queue labels contain three Priority, three Standard, and two Audit
  entries. Exchanging two equal labels does not create a new sequence.

  \[
    \frac{n!}{n_1!n_2!\cdots n_k!}
    =\frac{8!}{3!3!2!}=560
  \]

  \begin{alertblock}{The naive count overstates the sample space}
    \(8!=40{,}320\) treats repeated labels as distinguishable. Every distinct
    sequence appears \(3!3!2!=72\) times in that count.
  \end{alertblock}

  \textbf{Decision test:} Which exchanges change the final outcome?
\end{frame}

\begin{frame}[fragile]{Step 1: Formula and Enumeration Agree}
  \begin{columns}[T]
    \begin{column}{0.42\textwidth}
      \small
      \begin{lstlisting}
labels = (
    ["Priority"] * 3
    + ["Standard"] * 3
    + ["Audit"] * 2
)
distinct = set(permutations(labels))
      \end{lstlisting}

      \begin{block}{Verified output}
        Naive count: 40,320\\
        Formula: 560\\
        Enumerated sequences: 560\\
        Duplicate factor: 72
      \end{block}
    \end{column}
    \begin{column}{0.58\textwidth}
      \centering
      \includegraphics[width=\textwidth]{../figures/allocations_01_repeated_permutations.png}
    \end{column}
  \end{columns}
\end{frame}

\begin{frame}{Step 2: The Same Coefficient Allocates Labeled Requests}
  \begin{columns}[T]
    \begin{column}{0.40\textwidth}
      \small
      Allocate eight labeled requests to three labeled queues with sizes
      3, 3, and 2:
      \[
        \binom{8}{3,3,2}=\frac{8!}{3!3!2!}=560.
      \]

      Across all 560 allocations, R01 appears in:
      \begin{itemize}
        \item Priority: 210;
        \item Standard: 210;
        \item Audit: 140.
      \end{itemize}

      The 3:3:2 split follows the fixed queue sizes.
    \end{column}
    \begin{column}{0.60\textwidth}
      \centering
      \includegraphics[width=\textwidth]{../figures/allocations_02_multinomial_categories.png}
    \end{column}
  \end{columns}
\end{frame}

\begin{frame}{Step 3: Counting Occupancy Profiles}
  \small
  Distribute \(m\) identical units among \(r\) labeled categories. A profile
  \((x_1,\ldots,x_r)\) satisfies
  \[
    x_1+\cdots+x_r=m, \qquad x_i\geq0.
  \]

  The number of profiles is
  \[
    \binom{m+r-1}{r-1}.
  \]

  For eight units and four labeled service centers:
  \[
    \binom{8+4-1}{4-1}=\binom{11}{3}=165.
  \]

  \begin{exampleblock}{Verified computational evidence}
    Exhaustive Python enumeration also returns 165 ordered nonnegative
    occupancy profiles.
  \end{exampleblock}
\end{frame}

\begin{frame}{The Limit: A Possible Profile Is Not Necessarily Equiprobable}
  \begin{columns}[T]
    \begin{column}{0.37\textwidth}
      \small
      Treating all 165 profiles as equally likely gives
      \[
        \frac{1}{165}=0.00606061.
      \]

      Uniform independent routing instead gives
      \[
        P(x)=\frac{8!}{x_1!\cdots x_4!}
        \left(\frac14\right)^8.
      \]

      Seed-42 estimates from 500,000 simulations closely track these unequal
      exact probabilities.
    \end{column}
    \begin{column}{0.63\textwidth}
      \centering
      \includegraphics[width=\textwidth]{../figures/allocations_03_star_bars_probability_limit.png}
    \end{column}
  \end{columns}
\end{frame}

\begin{frame}{Concept Check}
  \begin{enumerate}
    \item How many sequences contain three P labels, two S labels, and two A
      labels?
    \item Why does the same coefficient count allocations of seven labeled
      requests to labeled queues of sizes 3, 2, and 2?
    \item How many ways can eight identical buffer units be distributed among
      three labeled centers?
    \item Are those occupancy profiles equally likely when eight labeled
      requests route independently? Explain.
  \end{enumerate}

  \vspace{0.4em}
  \begin{block}{Connection to Lesson 08}
    The next lesson formalizes probability rules and contingency tables. A
    count becomes a probability only after the probability model is justified.
  \end{block}

  \vfill
  \scriptsize Probability foundation: OpenStax,
  \textit{Introductory Business Statistics 2e}, Chapter 3, Section 3.1.
  Combinatorial formulas are a course-developed extension. CC BY-NC-SA 4.0.
\end{frame}

\appendix



\end{document}
